%=====================================================================
% Self-contained report. No external document is required to read it.
% Build:  cd tables/analysis && latexmk -pdf report_las_sto3g.tex
%=====================================================================
\documentclass[11pt]{article}
\usepackage[margin=2.5cm]{geometry}
\usepackage{amsmath,amssymb}
\usepackage{booktabs}
\usepackage{enumitem}
\usepackage{xcolor}
\usepackage[colorlinks=true,linkcolor=blue!55!black,
            urlcolor=blue!55!black]{hyperref}

\newcommand{\Ht}{\widetilde H}
\newcommand{\Uh}{\widehat U}
\newcommand{\FF}{\mathbb F_2}
\newcommand{\Ecal}{\mathcal E}
\newcommand{\Qcal}{\mathcal Q}
\newcommand{\Ccal}{\mathcal C}
\newcommand{\Gcal}{\mathcal G}
\newcommand{\Bcal}{\mathcal B}
\newcommand{\Hcal}{\mathcal H}
\newcommand{\Scal}{\mathcal S}
\newcommand{\Zb}{\bar Z}
\newcommand{\Edec}{E_{\mathrm{dec}}}
\newcommand{\Dmax}{D_{\max}}
\newcommand{\rsp}{r_{\mathrm{sp}}}
\usepackage{longtable}
\usepackage{amsthm}
\newtheorem{proposition}{Proposition}

\title{Admissibility conditions for a noncommutativity-ranked search for\\
local approximate parity symmetries, with an exhaustive control:\\
a numerical assessment on STO-3G $\mathrm{H_2O}$ and $\mathrm{N_2}$}
\author{}
\date{\today}

\begin{document}
\maketitle

\begin{abstract}
\begin{sloppypar}
We assess a search procedure that selects a fixed number $M$ of local
approximate parity symmetries (LASs) --- diagonal Pauli-$Z$ products that
almost commute with an electronic Hamiltonian --- by ranking candidates on a
state-specific noncommutativity (NC) score, quotienting out the exactly
conserved parities, and measuring the resulting sector structure through a
coupled dimension $K$. Three preconditions of the procedure are identified and
tested. First, the exact generators must commute with the Hamiltonian
\emph{at the geometry in use}; we show that storing point-group generators as
fixed orbital-index sets violates this along a dissociation scan, introduce the
retained reference weight $W$ as a run-time test, and find $W=0$ at every
$\mathrm{N_2}$ geometry of a prior calculation. Deriving the supports per
geometry from the Hartree--Fock (HF) irreducible-representation labels restores
$W=1$ at all $126$ points and with it the convergence of the $K$ criterion.
Second, the orbital rotation must map the exact set to diagonal Pauli products;
unrestricted $\mathrm{SO}(n)$ rotation does not, and we measure the resulting
leakage. Third, the LAS budget must satisfy $M\le n-\rsp$ on the spatial
subregister; we show that a prior configuration exceeded it by one row, and
correct it. An exhaustive control establishes that the polynomial current-frame
candidate library attains the exhaustive optimum, and we prove this is
guaranteed: the cost is constant on cosets of the exact group, so the candidate
classes form a linear matroid on which greedy ranking is optimal at every
budget. Two campaigns of $126$ points each, at $M=n-\rsp$ and $M=n-\rsp-1$,
find the two cost-ranked rules identical in cost, span and $K$ at all $84$
paired points, while a quota-constrained rule costs strictly more to select
yet never needs a larger coupled space.
\end{sloppypar}
\end{abstract}

%=====================================================================
\section{Introduction}
\label{sec:intro}

Exact Abelian symmetries of an electronic Hamiltonian that are diagonal in an
occupation basis --- particle-number parities, point-group parities --- can be
mapped to single qubits by a Clifford circuit and removed exactly. Operators
that \emph{almost} commute with the Hamiltonian admit no such exact treatment,
but fixing their eigenvalues defines a block decomposition whose diagonal
blocks are cheaper to diagonalise than the parent problem, at the cost of
discarding inter-block couplings. The procedure assessed here selects a fixed
number $M$ of such local approximate parity symmetries (LASs) by ranking
candidates on a state-specific noncommutativity score, and quantifies the
quality of the resulting decomposition by the number $K$ of block eigenstates
that must be recoupled to recover the target energy to a stated tolerance.

The procedure is stated in full in Sec.~\ref{sec:method}; no other document is
required. This report addresses four questions:

\begin{description}[leftmargin=2.6em,style=nextline,itemsep=2pt]
\item[Q1] Under what conditions is the exact-parity quotienting numerically
  valid, and can a violation be detected at run time?
\item[Q2] Is unrestricted $\mathrm{SO}(n)$ orbital rotation admissible when the
  exact set contains point-group parities?
\item[Q3] Does the polynomial candidate library, built from the current binary
  frame, reproduce an exhaustive scan over all parity classes?
\item[Q4] Does the choice of selection rule --- a quota-constrained greedy
  scan, an iterative re-ranked scan over the polynomial library, or an
  exhaustive scan --- change the pool selected or the coupled dimension $K$?
\end{description}

\noindent
Q1--Q4 are answered from the data reported here; Q4 is separated into a
question about the selected pool, which is settled, and one about $K$, which is
settled only in part and is reported as an observation (Sec.~\ref{sec:q4}). A
prior calculation on these systems is superseded: its searches never satisfied
the convergence criterion, so no coupled dimension was determined there and
none is quoted.

Acronyms used throughout: LAS (local approximate parity symmetry), NC
(noncommutativity), FCI (full configuration interaction), HF (Hartree--Fock),
RHF (restricted HF), MO (molecular orbital), JW (Jordan--Wigner), RREF (reduced
row-echelon form), PES (potential energy surface).

%=====================================================================
\section{Method}
\label{sec:method}

\subsection{Registers, Hamiltonian, orbital rotation}

Let $n$ be the number of spatial orbitals and $N=2n$ the number of occupation
qubits under the JW mapping with interleaved ordering: qubit $2p$ carries
spin-orbital $p\alpha$ and qubit $2p+1$ carries $p\beta$, $p=0,\dots,n-1$.
For an orbital rotation $U\in\mathrm{SO}(n)$ with number-conserving many-body
lift $\Uh$, the rotated Hamiltonian is
\begin{equation}
\Ht(U)=\Uh^{\dagger}H\Uh .
\label{eq:Htilde}
\end{equation}
$U$ is parameterised exponentially, not as a product of plane rotations:
\begin{equation}
U=\exp A(x),\qquad
A(x)_{ij}=x_k,\quad A(x)_{ji}=-x_k,\quad A(x)_{ii}=0,
\label{eq:expparam}
\end{equation}
where $k$ indexes an ordered list of pairs $(i,j)$ with $i<j$. The list is
taken in row-major upper-triangular order,
$(0,1),(0,2),\dots,(0,n-1),(1,2),\dots,(n-2,n-1)$, and comprises either all
$\binom{n}{2}$ pairs (\emph{unrestricted} packing) or only those pairs whose
two orbitals carry the same point-group irreducible representation
(\emph{irrep-block} packing), the surviving pairs retaining that order. Because
the generators $E_{ij}-E_{ji}$ do not commute, Eq.~\eqref{eq:expparam} is not
equal to the sequential product of Givens rotations through the same angles;
the two agree only to first order in $x$. The parameter vector $x$ recorded in
the artifacts of Appendix~\ref{sec:artifacts} is to be read through
Eq.~\eqref{eq:expparam}.
The rotation acts on the MO coefficient matrix as
\begin{equation}
C_{\mathrm{MO}}\;\longmapsto\;C_{\mathrm{MO}}\,U^{\mathsf T}.
\label{eq:convention}
\end{equation}

\subsection{Parity rows, exact group, quotient}

A binary row $g=(g_1,\dots,g_N)\in\FF^{N}$ labels the diagonal Pauli operator
\begin{equation}
Z(g)=\prod_{p=1}^{N}Z_p^{\,g_p},
\label{eq:Zg}
\end{equation}
and multiplication of such operators is addition of rows over $\FF$. Suppose
$r$ independent exact parities are known,
\begin{equation}
E_a=Z(e_a),\qquad [\Ht(U),E_a]=0,\qquad a=1,\dots,r,
\label{eq:exact}
\end{equation}
and set
\begin{equation}
\Ecal=\operatorname{span}_{\FF}\{e_1,\dots,e_r\}\subseteq\FF^{N},
\qquad
\Qcal=\FF^{N}/\Ecal\;\cong\;\FF^{\,N-r}.
\label{eq:quotient}
\end{equation}
LASs are classes in $\Qcal$, so the budget is bounded by
\begin{equation}
0\le M\le N-r .
\label{eq:budget}
\end{equation}
A candidate $g$ is accepted into the pool $\Gcal$ only if it increases the
rank of the stacked matrix,
\begin{equation}
\operatorname{rank}\begin{pmatrix}e_1;\dots;e_r;\,g_1;\dots;g_m;\,g\end{pmatrix}
>
\operatorname{rank}\begin{pmatrix}e_1;\dots;e_r;\,g_1;\dots;g_m\end{pmatrix},
\label{eq:ranktest}
\end{equation}
which in one test excludes exact rows, rows dependent on already accepted
LASs, and exact-dressed duplicates $g'=g+e$ with $e\in\Ecal$. The last
exclusion is legitimate because the cost is constant on the cosets of $\Ecal$
(Proposition~\ref{prop:coset}): $g$ and $g'$ define the same partition of that
sector.

\subsection{Cost functions and the reference state}

The target state $\Psi_{\mathrm{target}}$ is the exact FCI ground state of
$\Ht(U)$ in the fixed particle-number, spin-projection and exact-parity sector;
$E_{\mathrm{target}}$ is its energy. For a normalised reference
$|\Psi\rangle$, the NC score of a candidate is
\begin{equation}
c^{\mathrm{NC}}_{\Psi,U}(g)=\bigl\|\,[\Ht(U),Z(g)]\,|\Psi\rangle\,\bigr\|^{2},
\label{eq:nc}
\end{equation}
lower being preferred, and the alternative objective is the parity variance
\begin{equation}
c^{\mathrm{var}}_{\Psi,U}(g)=1-\langle\Psi|Z(g)|\Psi\rangle^{2}.
\label{eq:var}
\end{equation}
Both vanish when $Z(g)$ is an exact symmetry with a definite eigenvalue on
$\Psi$. They are evaluated in the determinant basis at fixed
$(N_\alpha,N_\beta)$, where $Z(g)$ is diagonal with signs $z_i=\pm1$, so that
\begin{equation}
\bigl([\Ht,Z(g)]\Psi\bigr)_i=\bigl(\Ht(z\odot\Psi)\bigr)_i-z_i\bigl(\Ht\Psi\bigr)_i .
\label{eq:commutator}
\end{equation}
$\Ht\Psi$ is formed once, after which Eq.~\eqref{eq:nc} costs one Hamiltonian
application per candidate and Eq.~\eqref{eq:var} costs none. The commutator is
anti-Hermitian, so $[\Ht,Z(g)]|\Psi\rangle$ has a purely imaginary expectation
structure; only its norm enters, and all arithmetic here is real.
% Insert into report_las_sto3g.tex immediately AFTER the subsection
% "Cost functions and the reference state" (i.e. after Eq.~\eqref{eq:commutator}
% and its paragraph), and BEFORE "Candidate libraries and the ranked scan".
%
% It promotes the claim already made informally at the end of the subsection
% "Parity rows, exact group, quotient" -- that c(g+e)=c(g) -- to a stated and
% proved proposition, and then draws the consequence that governs Q3 and Q4.

\subsection{The cost is a class function on the quotient}
\label{sec:classfunction}

The exclusion of exact-dressed duplicates in Eq.~\eqref{eq:ranktest} was
justified above by the assertion $c(g+e)=c(g)$. That identity is not a
convenience: it determines what any cost-ranked selection rule is able to
choose, and it is therefore stated and proved here.

\begin{proposition}[Coset invariance of the cost]
\label{prop:coset}
Let $|\Psi\rangle$ lie in a fixed exact sector, so that $Z(e)|\Psi\rangle
=\sigma_e|\Psi\rangle$ with $\sigma_e\in\{+1,-1\}$ for every $e\in\Ecal$. Then
for all $g\in\FF^{N}$ and all $e\in\Ecal$,
\begin{equation}
c^{\mathrm{NC}}_{\Psi,U}(g+e)=c^{\mathrm{NC}}_{\Psi,U}(g),
\qquad
c^{\mathrm{var}}_{\Psi,U}(g+e)=c^{\mathrm{var}}_{\Psi,U}(g).
\label{eq:cosetinvariance}
\end{equation}
\end{proposition}

\begin{proof}
Addition of rows over $\FF$ is multiplication of the operators in
Eq.~\eqref{eq:Zg}, so $Z(g+e)=Z(g)Z(e)$. By Eq.~\eqref{eq:exact},
$[\Ht(U),Z(e)]=0$, hence
\begin{equation}
[\Ht(U),Z(g)Z(e)]
=[\Ht(U),Z(g)]\,Z(e)+Z(g)\,[\Ht(U),Z(e)]
=[\Ht(U),Z(g)]\,Z(e).
\label{eq:commexpand}
\end{equation}
Applying this to $|\Psi\rangle$ and using $Z(e)|\Psi\rangle=\sigma_e|\Psi\rangle$
gives $[\Ht,Z(g+e)]|\Psi\rangle=\sigma_e\,[\Ht,Z(g)]|\Psi\rangle$. Since
$\sigma_e^{2}=1$, the norms in Eq.~\eqref{eq:nc} coincide. For
Eq.~\eqref{eq:var}, $\langle\Psi|Z(g)Z(e)|\Psi\rangle
=\sigma_e\langle\Psi|Z(g)|\Psi\rangle$, and the expectation is squared.
\end{proof}

Both cost functions are therefore constant on the cosets $g+\Ecal$: they are
well defined on the quotient $\Qcal$ of Eq.~\eqref{eq:quotient}, not merely on
$\FF^{N}$. Write $\bar c(\bar g)$ for the induced function, $\bar g$ denoting
the class of $g$.

Two structural consequences follow, and together they account for the
selection-rule results of Sections~\ref{sec:q3} and~\ref{sec:q4}.

\paragraph{(0) The ground set is the spatial subregister, not $\FF^{N}$.}
\begin{sloppypar}
Candidates are generated as products of the spin-symmetric pairs
$\Zb_p=Z_{p\alpha}Z_{p\beta}$, i.e.\ as images of the embedding $\iota$ of
Eq.~\eqref{eq:iota}, so the ground set of the matroid below is
$\iota(\FF^{n})\subset\FF^{N}$ and not $\FF^{N}$ itself. This is an
implementation restriction, not a requirement of the procedure: the candidate
generator, the GF(2) rank test and the cost evaluation all act on rows of
length $n$, and spin-resolved rows such as $Z_{p\alpha}$ alone are never
formed. Its consequences are that the accessible budget is bounded by
$n-\rsp$ rather than $N-r$ (Eq.~\eqref{eq:budget-sp}), and that the optima
reported here are optima over $\iota(\FF^{n})$ only; a spin-resolved search
could in principle reach a lower selection cost at the same $M$. Removing the
restriction requires the candidate library, the rank test and the cost to be
carried in the length-$N$ register throughout --- a change of representation
rather than of algorithm --- and it has not been done here. All statements in
Sections~\ref{sec:q3}, \ref{sec:q4} and~\ref{sec:budget} are therefore to be
read as holding within the spin-symmetric ground set.
\end{sloppypar}

\paragraph{(i) The pool is an independent set of a linear matroid.}
The acceptance test of Eq.~\eqref{eq:ranktest} admits $g$ exactly when
$\bar g$ is linearly independent of the classes already accepted. The
admissible pools of size $M$ are therefore precisely the independent sets of
size $M$ of the linear matroid on $\Qcal$, and the budget bound
Eq.~\eqref{eq:budget} is the statement that the rank of that matroid is
$\dim\Qcal$.

\paragraph{(ii) Greedy ranking is optimal at every budget.}
For a matroid with non-negative element weights, the greedy algorithm returns a
minimum-weight independent set of size $M$ for \emph{every} $M$, not only at
full rank. Since the weight $\bar c$ is well defined on $\Qcal$ by
Proposition~\ref{prop:coset}, any two selection rules that both attain the
greedy optimum select independent sets of equal total weight, and hence
subspaces $\operatorname{span}\{\bar g_1,\dots,\bar g_M\}\subseteq\Qcal$ that
are optimal in the same sense. Minimum-weight independent sets need not be
unique; when they are not, the rules may return different rows while spanning
the same subspace.

\paragraph{Gauge freedom.}
A rule reports rows $g_j\in\FF^{N}$, but only the classes $\bar g_j$ carry
information: by Proposition~\ref{prop:coset} the cost is unchanged, and by
Eq.~\eqref{eq:ranktest} the accepted rank is unchanged, if any $g_j$ is
replaced by $g_j+e$. The group generated by $\Ecal$ together with the pool,
and therefore the sector partition, the block dimensions $\Dmax$ and the
coupled dimension $K$, depend only on
$\operatorname{span}\{\bar g_1,\dots,\bar g_M\}$. Differences between rules in
the reported rows are consequently \emph{not} evidence of a different
selection; only a difference in the spanned subspace is. This distinction is
applied throughout Section~\ref{sec:q4}, where rules are compared by span
rather than by row equality.

\paragraph{A degenerate budget.}
If $M=\dim\Qcal$ then any $M$ independent classes form a basis of $\Qcal$, so
$\operatorname{span}\{\bar g_1,\dots,\bar g_M\}=\Qcal$ for \emph{every}
admissible pool, and
$\operatorname{span}(\Ecal\cup\Gcal)=\FF^{N}$ irrespective of the rule. At that
budget the partition is rule-independent by construction and carries no
discriminating information; a rule can influence $K$ only indirectly, through
the orbital objective of Eq.~\eqref{eq:objective}. This degeneracy is the one
noted in Section~\ref{sec:matroid} for the span test; the statement here is
its general form, and it applies to $\Dmax$ and the block count as well.
Comparisons at $M=\dim\Qcal$ must therefore be read with the degeneracy in
mind, and Section~\ref{sec:budget} repeats the campaign at $M=\dim\Qcal-1$,
where the spanned subspace is proper and the degeneracy is absent.



\subsection{Candidate libraries and the ranked scan}

At macroiteration $t$ the search maintains a full binary frame
$F^{(t)}=(f^{(t)}_1;\dots;f^{(t)}_N)\in\mathrm{GL}(N,\FF)$ containing a basis of
$\Ecal$, the current $M$ LAS rows, and $N-r-M$ auxiliary rows. The
\emph{polynomial} candidate library consists of the current-frame single rows
and their pairwise sums,
\begin{equation}
\begin{gathered}
\Ccal^{(t)}=\{f^{(t)}_j\}_{j\le N}\;\cup\;\{f^{(t)}_j+f^{(t)}_k\}_{j<k},\\[2pt]
|\Ccal^{(t)}|=\tfrac{N(N+1)}{2},
\end{gathered}
\label{eq:Cpoly}
\end{equation}
whereas the \emph{exhaustive} library is the complete quotient,
\begin{equation}
\Ccal_{\mathrm{full}}=\{\bar g\in\FF^{\,N-r}:\bar g\neq0\},
\quad |\Ccal_{\mathrm{full}}|=2^{\,N-r}-1 .
\label{eq:Cfull}
\end{equation}
In either mode the library is sorted by ascending cost and scanned in full,
applying Eq.~\eqref{eq:ranktest}. The first $M$ accepted rows are the LAS pool
$\Gcal^{(t+1)}$. \emph{The scan does not stop there}: it continues until
$N-r$ rows have been accepted, forming a complete ordered quotient basis
$\Bcal^{(t+1)}$. The reason is frame construction, not the state route: the
trailing $N-r-M$ rows become the auxiliary axes of the next frame
$F^{(t+1)}$, and hence define the next current-frame neighbourhood
Eq.~\eqref{eq:Cpoly}. Continuing the same ranked scan removes what would
otherwise be an arbitrary per-iteration choice of completion rows. Ties are
broken deterministically by ascending (cost, frame index).

\subsection{Orbital objective}

The reference is reconstructed at every trial rotation --- here as the exact
FCI ground state of $\Ht(U)$ restricted to the target sector --- and the
orbital objective is the generator sum
\begin{equation}
C_A(U;\Gcal)=\sum_{i=1}^{M}c_{\Psi(U),U}\bigl(g_i\bigr),
\qquad
U^{(t)}=\arg\min_U C_A\bigl(U;\Gcal^{(t)}\bigr).
\label{eq:objective}
\end{equation}
Equation~\eqref{eq:objective} is a sum over a chosen generator list and is
therefore basis dependent: two lists spanning the same quotient subspace may
give different values. It is not an invariant of the LAS partition.

\subsection{Clifford reduction, decoupled energy, coupled dimension}

A Clifford circuit maps the $r$ exact rows and the $M$ LAS rows to distinct
single-qubit $Z$ operators. Fixing the $r$ exact labels to the physical values
$\boldsymbol e$ is an exact reduction; fixing the $M$ LAS labels to
$\boldsymbol s$ additionally discards couplings between approximate sectors and
is therefore an approximation. The resulting block acts on
$N_{\mathrm{tail}}=N-r-M$ qubits. $\Edec$ denotes the lowest eigenvalue over
the retained blocks.

To quantify the approximation, the block eigenstates $|\phi_j\rangle$ of the
fixed exact sector are ordered by decreasing target weight
\begin{equation}
w_j=\bigl|\langle\phi_j|\Psi_{\mathrm{target}}\rangle\bigr|^{2},
\label{eq:wj}
\end{equation}
the subspace matrices $H^{(K)}_{ij}=\langle\phi_i|\Ht(U)|\phi_j\rangle$ are
formed, and the coupled dimension $K$ is the smallest integer satisfying
\begin{equation}
\lambda_{\min}\bigl(H^{(K)}\bigr)-E_{\mathrm{target}}<\varepsilon_E,
\qquad \varepsilon_E=1.6\times10^{-3}\ \mathrm{Ha}.
\label{eq:K}
\end{equation}
$K$ is a retrospective benchmark requiring knowledge of the target state; it is
not itself a deployable selection rule.

\subsection{The spatial subregister}
\label{sec:subregister}

The implementation examined here represents LAS candidates exclusively as
\emph{spin-symmetric} products
\begin{equation}
\Zb_p=Z_{p\alpha}Z_{p\beta},
\end{equation}
that is, inside the image of the linear embedding
$\iota:\FF^{\,n}\hookrightarrow\FF^{\,N}$ defined by
\begin{equation}
\bigl(\iota(x)\bigr)_{2p}=\bigl(\iota(x)\bigr)_{2p+1}=x_p,
\qquad p=0,\dots,n-1,
\label{eq:iota}
\end{equation}
i.e.\ each spatial bit is duplicated onto both spin qubits of that orbital.
This is a hard limitation of the implementation, not a run-time option: the
binary frame is constructed over $n$ bits and every generator is emitted as a
$Z$ product over spatial orbital indices, so no code path can return a
spin-resolved row such as $Z_{p\alpha}$ alone.

Writing $\Ecal_{\mathrm{sp}}=\iota^{-1}(\Ecal)$ and
$\rsp=\dim\Ecal_{\mathrm{sp}}$, the admissible budget inside the subregister is
\begin{equation}
M\le n-\rsp ,
\label{eq:budget-sp}
\end{equation}
in general stronger than Eq.~\eqref{eq:budget}. Note that total particle parity
\begin{equation}
\mathbf 1=\prod_{p}\Zb_p=(-1)^{\hat N}=P_\alpha\oplus P_\beta
\label{eq:allones}
\end{equation}
lies in $\Ecal_{\mathrm{sp}}$ whenever the spin-resolved number parities
$P_\alpha,P_\beta$ belong to $\Ecal$.

\subsection{Retained reference weight}
\label{sec:W}

We introduce the diagnostic
\begin{equation}
W=\sum_{j}w_j ,
\label{eq:W}
\end{equation}
summed over all block eigenstates retained after exact tapering. If the pinned
operators are genuine symmetries of $\Ht(U)$ and the pinned sector is the one
containing $\Psi_{\mathrm{target}}$, then the retained states span that sector
and $W=1$ identically. If $W<1$, a fraction $1-W$ of the target state lies
outside the retained space and is unreachable by any $K$, so Eq.~\eqref{eq:K}
cannot be satisfied however the pool is chosen. $W$ is therefore a sufficient
run-time test of the validity of the exact quotienting, at the cost of one
overlap sum. We adopt $1-W<10^{-4}$ as a precondition.

%=====================================================================
\section{Computational details}
\label{sec:details}

\paragraph{Systems and geometries.}
$\mathrm{H_2O}$/STO-3G: $n=7$, $N=14$, $N_\alpha=N_\beta=5$, $C_{2v}$, symmetric
stretch at fixed $\angle\mathrm{HOH}=104.5^{\circ}$, ten geometries uniformly
spaced from $R_{\mathrm{OH}}=0.958$\,\AA\ to $2.5$\,\AA\ in steps of
$(2.5-0.958)/9=0.171\overline{3}$\,\AA. Exact set $r=4$:
$P_\alpha,P_\beta,Q_{B1},Q_{B2}$.
$\mathrm{N_2}$/STO-3G: $n=10$, $N=20$, $N_\alpha=N_\beta=7$, $D_{2h}$, eleven
geometries $R_{\mathrm{NN}}=1.2(0.1)2.2$\,\AA. Exact set $r=5$:
$P_\alpha,P_\beta,Q_{\pi x},Q_{\pi y},Q_u$. The point-group parity $Q_\Gamma$ is
the product of $\Zb_p$ over the orbitals belonging to the irreducible
representations collected in $\Gamma$.

\paragraph{Selection rules and grid size.}
Two rules are implemented. The \emph{greedy quota} selects
$n_{\mathrm{singles}}$ single-$\Zb$ and $n_{\mathrm{quartets}}$ pair rows once,
at the initial frame, in one pass (a variant additionally forbids two selected
rows to share an orbital). The \emph{iterative} rule performs the ranked scan
of Sec.~\ref{sec:method} anew at every macroiteration. With three rule variants
and two objectives at $10+11=21$ geometries, the grid comprises $126$ points.

\paragraph{Initialisation and determinism.}
The orbital basis is the canonical RHF MOs of the geometry. The initial binary
frame is $F^{(0)}=I$, meaning the frame rows are the unit vectors of the
spatial parity register, i.e.\ the single-orbital parities $\Zb_p$. The initial
rotation is the identity, $x^{(0)}=0$. The orbital optimiser is L-BFGS-B with
finite-difference gradients, step tolerance $10^{-5}$, gradient tolerance
$10^{-6}$, at most $100$ iterations. At most $20$ macroiterations are performed;
the search stops when the LAS quotient span is unchanged for two consecutive
macroiterations. At most $500$ eigenstates are computed per block. No random
number generator enters the selection or the optimisation: given the geometry
and the settings, the calculation is deterministic.

\paragraph{Exact sector.}
The pinned exact sector is the one containing the aufbau reference determinant,
obtained by applying the Clifford map to its occupation string. It is not
chosen by sector size.

\paragraph{Software.}
Integrals, RHF, orbital symmetry labels and FCI from PySCF; qubit operators and
the Clifford construction from OpenFermion and an in-house package. Exact
version identifiers are recorded in the artifact manifest of
Appendix~\ref{sec:artifacts}.

%=====================================================================
\section{Q1: validity of the exact quotienting}
\label{sec:q1}

\subsection{Failure mode}
Equation~\eqref{eq:exact} demands $[\Ht(U),E_a]=0$. In the prior calculation the
point-group generator supports were stored as fixed orbital index sets, derived
once at a single reference geometry. MO orderings cross along both dissociation
scans, so at other geometries those indices designate orbitals of different
irreducible representations, and the corresponding $Z$ products do not commute
with $\Ht(U)$. The exact tapering then projects onto a sector that does not
contain the target state, and Eq.~\eqref{eq:K} becomes unsatisfiable.

\begin{table}[htbp]
\centering
\caption{Retained reference weight $W$ in the prior calculation (unrestricted
packing, iterative rule, NC objective), against agreement of the stored
generator supports with the HF irreducible-representation labels at each
geometry. Where the supports disagree, Eq.~\eqref{eq:K} is not satisfied and no
coupled dimension is determined.}
\label{tab:q1}
\begin{tabular}{lccc}
\toprule
$R$ (\AA) & supports agree & $W$ & Eq.~\eqref{eq:K} satisfied \\
\midrule
\multicolumn{4}{l}{\emph{$\mathrm{H_2O}$}}\\
$0.958,\ 1.1293,\ 1.3007$ & yes & $1.0000$ & yes \\
1.4720 & no & 0.9255 & no \\
1.6433 & no & 0.8855 & no \\
1.8147 & no & 0.8424 & no \\
1.9860 & no & 0.8054 & no \\
2.1573 & no & 0.7801 & no \\
2.3287 & no & 0.7655 & no \\
2.5000 & no & 0.0242 & no \\
\midrule
\multicolumn{4}{l}{\emph{$\mathrm{N_2}$, all eleven geometries}}\\
$1.2$--$2.2$ & no & $\mathbf{0.0000}$ & no \\
\bottomrule
\end{tabular}
\end{table}

Two features identify the mechanism unambiguously. For $\mathrm{H_2O}$ the
boundary between $W=1$ and $W<1$ coincides exactly with the $3a_1/1b_1$ level
crossing at $R=1.472$\,\AA. For $\mathrm{N_2}$ the stored supports are already
wrong at $R=1.2$\,\AA\ --- the $3\sigma_g$ orbital occupies index $4$, so the
correct supports are $Q_{\pi x}=\{6,7\}$ and $Q_u=\{1,3,5,6,9\}$ --- and $W=0$
throughout: the retained space is orthogonal to the target state.

\subsection{Correction}
Deriving the supports at each geometry from the point-group labels of the
converged HF orbitals, and pinning the sector containing the reference
determinant, gives $W=1$ and a correctly counted exact block at all $126$
points, with Eq.~\eqref{eq:K} satisfied at every one. Convergence is restored
throughout. The resulting coupled dimensions are obtained at a budget that
violates Eq.~\eqref{eq:budget-sp} (Sec.~\ref{sec:q4}) and are therefore not
interpreted here.

One structural observation is independent of the budget: the largest retained
block is
\begin{equation}
\Dmax=\binom{n}{N_\alpha}=21\ (\mathrm{H_2O}),\qquad 120\ (\mathrm{N_2}),
\end{equation}
at every geometry, rule, objective and packing. This is the seniority-zero
block, in which every spatial orbital is doubly occupied or empty: when the
span of $\Ecal\cup\Gcal$ exhausts the register, each block is a single
occupation-parity pattern and the closed-shell pattern is always the largest.

%=====================================================================
\section{Q2: admissible orbital rotations}
\label{sec:q2}

A point-group parity $Q_\Gamma$ is defined by a set of orbital indices. Under
an orbital rotation the operator is carried to $\Uh^{\dagger}Q_\Gamma\Uh$. If
$U$ is block diagonal in the irreducible representations, this is again a
product of the same $\Zb_p$ and remains a diagonal Pauli operator, so the
Clifford construction of Sec.~\ref{sec:method} applies unchanged. If $U$ mixes
orbitals of different irreducible representations, the transformed operator is
no longer diagonal in the occupation basis, is not a Pauli $Z$ product, and
cannot be mapped to a single qubit by the Clifford circuit used. Spin-resolved
number parities are unaffected, since any spin-conserving rotation preserves
$N_\alpha$ and $N_\beta$.

The consequence is measurable. At a fixed geometry the exact generators are
identical across selection rules, and only the optimised $U$ differs; the
leakage in Table~\ref{tab:q2} is therefore attributable to the rotation alone.
It grows monotonically with bond length, consistent with the optimiser
exploiting inter-irrep mixing more strongly as correlation increases.

\begin{table}[htbp]
\centering
\caption{Exact-symmetry leakage $1-W$ under unrestricted $\mathrm{SO}(n)$
rotation ($\mathrm{H_2O}$, greedy disjoint quota, variance objective; generator
supports correct at these geometries). The irrep-block packing gives $W=1$ at
all of them.}
\label{tab:q2}
\begin{tabular}{lcccccc}
\toprule
$R$ (\AA) & 1.1293 & 1.3007 & 1.4720 & 1.6433 & 1.8147 & 1.9860 \\
\midrule
$W$   & 0.9763 & 0.9522 & 0.9169 & 0.8667 & 0.8073 & 0.7154 \\
$1-W$ & 2.4\%  & 4.8\%  & 8.3\%  & 13.3\% & 19.3\% & 28.5\% \\
\bottomrule
\end{tabular}
\end{table}

\paragraph{Relation to Table~\ref{tab:q1}.}
Tables~\ref{tab:q1} and~\ref{tab:q2} both report unrestricted-packing runs and
must not be compared row by row. $W$ depends on the optimised rotation, which
depends on the selection rule and objective; Table~\ref{tab:q1} lists the
iterative rule with the NC objective, Table~\ref{tab:q2} the greedy disjoint
quota with the variance objective. The two tables isolate different causes:
Table~\ref{tab:q1} the incorrect generator supports, Table~\ref{tab:q2} the
rotation, at geometries where the supports are correct.

Where the unrestricted packing does retain the reference it reproduces the
irrep-block result: of the $119$ points for which both packings completed
successfully, $118$ agree exactly in $K$ and one differs by a single unit.

%=====================================================================
\section{Q3: exhaustive control}
\label{sec:q3}

\subsection{Construction}
At each point's converged orbitals, every class of $\Ccal_{\mathrm{full}}$ in
the spatial subregister was scored with the objective that run had used, using
Eq.~\eqref{eq:commutator}; the ranked scan with the test
Eq.~\eqref{eq:ranktest} was repeated on the complete ranking; and the outcome
was compared with the polynomial pool. The control was run with the same exact
set as the calculation under test, giving pool sizes $2^{\,n-\rsp}-1=31$
($\mathrm{H_2O}$) and $127$ ($\mathrm{N_2}$); with $\mathbf 1$ of
Eq.~\eqref{eq:allones} restored to $\Ecal$ (Sec.~\ref{sec:q4}) the
corresponding sizes would be $15$ and $63$.

A correctness gate re-scores the polynomial pool's own rows and compares with
the values recorded by the production code. $121$ of $126$ points reproduce
them to within $10^{-4}$ relative. The five exceptions are $\mathrm{H_2O}$ at
$R\ge2.33$\,\AA, where the FCI gap is $E_1-E_0=4.5\times10^{-4}$\,Ha at
$R=2.5$\,\AA: the reference is then not determined by energy alone and
independent solvers may return different vectors from the near-degenerate
manifold. Those points are excluded.

Of the remaining points, the total-objective comparison requires both pools to
be scored at the \emph{same} orbitals. This holds for the $41$ surviving
iterative points. It does not hold for the greedy points, whose recorded costs
were evaluated before orbital optimisation (Sec.~\ref{sec:q4}, item iii), so
they are excluded from this comparison.

\subsection{Result}
Over those $41$ points,
\begin{equation}
\frac{\sum_i c\bigl(g_i^{\,\mathrm{exhaustive}}\bigr)}
     {\sum_i c\bigl(g_i^{\,\mathrm{polynomial}}\bigr)}
\in\bigl[1-5\times10^{-12},\;1\bigr],
\end{equation}
i.e.\ the two totals agree to $10^{-9}$ at every one.

\subsection{Why this is a non-trivial test}
\label{sec:matroid}
Linear independence over $\FF$ defines a matroid on the set of parity classes.
Equation~\eqref{eq:ranktest} is precisely the greedy rank test for that
matroid, and greedy selection in ascending weight on a matroid returns a
minimum-weight independent set of size $M$ of the ground set it is. The exhaustive and
given, for every $M$ (Section~\ref{sec:classfunction}(ii)). polynomial scans are the same greedy algorithm on the same matroid, differing
only in ground set: $\Ccal_{\mathrm{full}}$ versus $\Ccal^{(t)}$. Equality of
the two totals therefore states that the current-frame neighbourhood
$\Ccal^{(t)}$ contains a minimum-weight basis of the full quotient matroid.

The span comparison is degenerate at this budget and the total-objective
test is not; both statements are made in general form in
Section~\ref{sec:classfunction}, and are not repeated here.

%=====================================================================
% Replaces the section currently titled
%   "Q4: obstructions to comparing the selection rules"
% Requires: \label{sec:classfunction} (Proposition~\ref{prop:coset}),
%           \label{sec:budget}, and the generated tables
%           gen_pool_equivalence.tex, gen_k_by_arm.tex.

\section{Q4: comparison of the selection rules}
\label{sec:q4}

\subsection{Statement of the question}

The question posed in the introduction conflates two claims that admit answers
of very different strength, and they are separated here.

\begin{description}
\item[Q4(i)] Does the polynomial candidate library of
  Eq.~\eqref{eq:Cpoly} select the same pool as an exhaustive search over
  the full candidate space?
\item[Q4(ii)] Does the choice of selection rule change the coupled dimension
  $K$?
\end{description}

Q4(i) is answered affirmatively, and Proposition~\ref{prop:coset} shows why the
answer is forced rather than incidental. Q4(ii) is answered only in part: a
difference is observed for one rule, but at a single orbital start it is not
separable from optimiser scatter, and it is reported as an observation.

\subsection{Design}

Three rules are compared at fixed budget, cost function, molecule and geometry:

\begin{description}
\item[exhaustive] greedy ranking over the full candidate space, i.e.\ the
  minimum-weight independent set of the matroid of
  Section~\ref{sec:classfunction}(i);
\item[iterative] greedy ranking restricted to the polynomial library
  $\Ccal^{(t)}$, re-formed at each macroiteration;
\item[quota] greedy ranking constrained to a prescribed split of single-orbital
  and quartet supports, $(n_{\mathrm{s}},n_{\mathrm{q}})=(2,\,2)$ for
  $\mathrm{H_2O}$ and $(3,\,3)$ for $\mathrm{N_2}$ at $M=\dim\Qcal$, and $(2,\,1)$ and $(3,\,2)$
  respectively at $M=\dim\Qcal-1$.
\end{description}

\paragraph{Choice of the quota split.}
The split $(n_{\mathrm{s}},n_{\mathrm{q}})$ is a prescribed input, not a
quantity the procedure derives, and it must satisfy
$n_{\mathrm{s}}+n_{\mathrm{q}}=M$. At the full budget the balanced splits
$(2,\,2)$ and $(3,\,3)$ were adopted. At the reduced budget the deficit of one was
taken from the quartets, giving $(2,\,1)$ and $(3,\,2)$, so that the number of
single-orbital rows is held fixed across the two budgets and the two campaigns
differ in one respect only. The complementary choices $(1,\,2)$ and $(2,\,3)$ were
not run. Consequently the results below characterise \emph{these} splits and
not quota selection in general: a different split could in principle come
closer to, or reach, the unconstrained optimum. This is a limitation of the
present campaign rather than a property of quota rules, and it is repeated in
Section~\ref{sec:q4limits}.

\paragraph{Sign convention.}
\begin{sloppypar}
Throughout this section
\begin{equation}
\Delta K \;=\; K_{\mathrm{quota}}-K_{\mathrm{exhaustive}},
\label{eq:dK}
\end{equation}
\noindent evaluated at fixed molecule, geometry, cost function and budget, so
that a \emph{negative} $\Delta K$ means the quota rule required a smaller coupled
space than the exhaustive rule. Comparisons between the iterative and
exhaustive rules are reported as agreement or disagreement rather than as a
signed difference, since they never disagree.
\end{sloppypar}

\paragraph{Counting of points.}
The two campaigns comprise $252$ runs: $21$ geometries ($10$ for
$\mathrm{H_2O}$, $11$ for $\mathrm{N_2}$) $\times$ $2$ cost functions $\times$
$3$ rules $\times$ $2$ budgets. A \emph{pairwise comparison} of two rules
therefore has $21\times2\times2=84$ points, and a statement restricted to one
molecule has $10\times2\times2=40$ points for $\mathrm{H_2O}$ and
$11\times2\times2=44$ for $\mathrm{N_2}$.

Each rule is followed by the same orbital optimisation of
Eq.~\eqref{eq:objective}, started from the identity, and by the same Clifford
reduction and coupled-dimension measurement, Eq.~\eqref{eq:K}. The grids are
ten $\mathrm{H_2O}$ geometries
($R=0.958$--$2.5$\,\AA) and eleven $\mathrm{N_2}$ geometries ($R=1.2$--$2.2$\,\AA),
each with both cost functions and all three rules, at two budgets: $126$
points per budget, $252$ in total. Every point satisfies the preconditions of
Section~\ref{sec:verification} --- retained reference weight $W=1$ to within
$2.2\times10^{-7}$, $n_{\mathrm{exact}}=r$, and $M$ at the recorded budget ---
so no point is excluded from what follows.

Because the rows reported by a rule are defined only up to $\Ecal$
(Section~\ref{sec:classfunction}), rules are compared by the spanned subspace
$\operatorname{span}\{\bar g_1,\dots,\bar g_M\}\subseteq\Qcal$ and by the total
selection cost $\sum_j \bar c(\bar g_j)$, never by row equality.

\subsection{Q4(i): the polynomial library attains the exhaustive optimum}

\begin{table}[htbp]
\centering
\begin{tabular}{llrrrr}
\hline
budget & cost & points & $\max|\Delta c|$ & span agree & $K$ agree\\
\hline
$M=n-\rsp$ & NC & 21 & $2.4\times10^{-14}$ & 21/21 & 21/21\\
$M=n-\rsp$ & variance & 21 & $3.0\times10^{-15}$ & 21/21 & 21/21\\
$M=n-\rsp-1$ & NC & 21 & $2.7\times10^{-14}$ & 21/21 & 21/21\\
$M=n-\rsp-1$ & variance & 21 & $2.2\times10^{-15}$ & 21/21 & 21/21\\
\hline
\end{tabular}

\caption{Agreement between the polynomial (iterative) and exhaustive rules.
``points'' counts geometries at the stated budget and cost function;
$\Delta c=\sum_j\bar c(\bar g_j)$ differs between the two rules;
``span agree'' counts points at which
$\operatorname{span}\{\bar g_j\}$ coincides; ``$K$ agree'' counts points at
which the coupled dimension coincides. Row sets never coincide.}
\label{tab:poolequiv}
\end{table}

Table~\ref{tab:poolequiv} reports agreement in three distinct senses, and the
distinction matters. Across both budgets and both cost functions the two rules
agree in \emph{total selection cost} to floating-point round-off, in the
\emph{spanned subspace} of $\Qcal$ at every point, and in the \emph{coupled
dimension} $K$ at every point. They do \emph{not} agree in the reported rows:
at no point of either campaign do the two rules return the same row set. Where
this section says the rules are ``identical'', it therefore means identical in
$\bar c$, in span, and in $K$ --- never in rows.

That combination is exactly what Section~\ref{sec:classfunction}(ii) requires:
minimum-weight independent sets of equal total weight need not coincide as
sets, and the residual freedom is the choice of coset representative, which by
Proposition~\ref{prop:coset} is unobservable. The agreement is therefore not
evidence that the polynomial library happens to be adequate; it is a
consequence of greedy optimality on the matroid, and the measurement confirms
that the implementation realises it.

The practical consequence is that the polynomial library is not an
approximation to the exhaustive search at these system sizes: it attains the
same optimum at polynomial rather than exponential candidate cost.

\subsection{Q4(ii): the effect of the selection rule on \texorpdfstring{$K$}{K}}

% h2o NC M=n-\rsp: lower 0, equal 10, higher 0, mean dK = +0.00
% h2o variance M=n-\rsp: lower 0, equal 10, higher 0, mean dK = +0.00
% n2 NC M=n-\rsp: lower 5, equal 4, higher 2, mean dK = -0.27
% n2 variance M=n-\rsp: lower 1, equal 10, higher 0, mean dK = -0.09
% h2o NC M=n-\rsp-1: lower 0, equal 10, higher 0, mean dK = +0.00
% h2o variance M=n-\rsp-1: lower 0, equal 10, higher 0, mean dK = +0.00
% n2 NC M=n-\rsp-1: lower 5, equal 5, higher 1, mean dK = -0.45
% n2 variance M=n-\rsp-1: lower 1, equal 10, higher 0, mean dK = -0.09
\begin{longtable}{lrrr}
\caption{Coupled dimension $K$ by selection rule, geometry, cost function and budget. Each block is headed by its molecule, cost function and budget. ``quota'' is the constrained rule with the splits of Section~\ref{sec:q4}; $\Delta K$ counts are in Table~\ref{tab:dk}.}
\label{tab:kbyarm}\\
\hline
$R$ (\AA) & exhaustive & iterative & quota\\
\hline
\endfirsthead
\multicolumn{4}{l}{\emph{Table \thetable\ continued}}\\
\hline
$R$ (\AA) & exhaustive & iterative & quota\\
\hline
\endhead
\hline
\endfoot
\multicolumn{4}{l}{}\\
\multicolumn{4}{l}{\emph{$\mathrm{H_2O}$, NC cost, $M=n-\rsp$}}\\
\hline
0.9580 & 4 & 4 & 4\\
1.1293 & 5 & 5 & 5\\
1.3007 & 7 & 7 & 7\\
1.4720 & 8 & 8 & 8\\
1.6433 & 8 & 8 & 8\\
1.8147 & 7 & 7 & 7\\
1.9860 & 5 & 5 & 5\\
2.1573 & 5 & 5 & 5\\
2.3287 & 4 & 4 & 4\\
2.5000 & 5 & 5 & 5\\
\multicolumn{4}{l}{}\\
\multicolumn{4}{l}{\emph{$\mathrm{H_2O}$, variance cost, $M=n-\rsp$}}\\
\hline
0.9580 & 3 & 3 & 3\\
1.1293 & 5 & 5 & 5\\
1.3007 & 6 & 6 & 6\\
1.4720 & 7 & 7 & 7\\
1.6433 & 7 & 7 & 7\\
1.8147 & 6 & 6 & 6\\
1.9860 & 5 & 5 & 5\\
2.1573 & 5 & 5 & 5\\
2.3287 & 4 & 4 & 4\\
2.5000 & 5 & 5 & 5\\
\multicolumn{4}{l}{}\\
\multicolumn{4}{l}{\emph{$\mathrm{N_2}$, NC cost, $M=n-\rsp$}}\\
\hline
1.2000 & 23 & 23 & 26\\
1.3000 & 26 & 26 & 25\\
1.4000 & 26 & 26 & 26\\
1.5000 & 26 & 26 & 27\\
1.6000 & 24 & 24 & 22\\
1.7000 & 23 & 23 & 21\\
1.8000 & 18 & 18 & 17\\
1.9000 & 14 & 14 & 13\\
2.0000 & 9 & 9 & 9\\
2.1000 & 7 & 7 & 7\\
2.2000 & 4 & 4 & 4\\
\multicolumn{4}{l}{}\\
\multicolumn{4}{l}{\emph{$\mathrm{N_2}$, variance cost, $M=n-\rsp$}}\\
\hline
1.2000 & 20 & 20 & 20\\
1.3000 & 22 & 22 & 21\\
1.4000 & 22 & 22 & 22\\
1.5000 & 23 & 23 & 23\\
1.6000 & 22 & 22 & 22\\
1.7000 & 21 & 21 & 21\\
1.8000 & 17 & 17 & 17\\
1.9000 & 13 & 13 & 13\\
2.0000 & 9 & 9 & 9\\
2.1000 & 7 & 7 & 7\\
2.2000 & 4 & 4 & 4\\
\multicolumn{4}{l}{}\\
\multicolumn{4}{l}{\emph{$\mathrm{H_2O}$, NC cost, $M=n-\rsp-1$}}\\
\hline
0.9580 & 1 & 1 & 1\\
1.1293 & 1 & 1 & 1\\
1.3007 & 1 & 1 & 1\\
1.4720 & 1 & 1 & 1\\
1.6433 & 1 & 1 & 1\\
1.8147 & 1 & 1 & 1\\
1.9860 & 1 & 1 & 1\\
2.1573 & 1 & 1 & 1\\
2.3287 & 1 & 1 & 1\\
2.5000 & 3 & 3 & 3\\
\multicolumn{4}{l}{}\\
\multicolumn{4}{l}{\emph{$\mathrm{H_2O}$, variance cost, $M=n-\rsp-1$}}\\
\hline
0.9580 & 1 & 1 & 1\\
1.1293 & 1 & 1 & 1\\
1.3007 & 1 & 1 & 1\\
1.4720 & 1 & 1 & 1\\
1.6433 & 1 & 1 & 1\\
1.8147 & 1 & 1 & 1\\
1.9860 & 1 & 1 & 1\\
2.1573 & 1 & 1 & 1\\
2.3287 & 1 & 1 & 1\\
2.5000 & 3 & 3 & 3\\
\multicolumn{4}{l}{}\\
\multicolumn{4}{l}{\emph{$\mathrm{N_2}$, NC cost, $M=n-\rsp-1$}}\\
\hline
1.2000 & 16 & 16 & 15\\
1.3000 & 18 & 18 & 17\\
1.4000 & 18 & 18 & 19\\
1.5000 & 18 & 18 & 18\\
1.6000 & 17 & 17 & 16\\
1.7000 & 25 & 25 & 23\\
1.8000 & 19 & 19 & 19\\
1.9000 & 15 & 15 & 14\\
2.0000 & 7 & 7 & 7\\
2.1000 & 4 & 4 & 4\\
2.2000 & 3 & 3 & 3\\
\multicolumn{4}{l}{}\\
\multicolumn{4}{l}{\emph{$\mathrm{N_2}$, variance cost, $M=n-\rsp-1$}}\\
\hline
1.2000 & 11 & 11 & 11\\
1.3000 & 15 & 15 & 14\\
1.4000 & 17 & 17 & 17\\
1.5000 & 16 & 16 & 16\\
1.6000 & 15 & 15 & 15\\
1.7000 & 15 & 15 & 15\\
1.8000 & 11 & 11 & 11\\
1.9000 & 7 & 7 & 7\\
2.0000 & 5 & 5 & 5\\
2.1000 & 4 & 4 & 4\\
2.2000 & 2 & 2 & 2\\
\end{longtable}


Three observations follow, in decreasing order of evidential strength.

\paragraph{Cost-ranked rules do not differ.}
The exhaustive and iterative columns of Table~\ref{tab:kbyarm} carry the same
$K$ at all $84$ points, in the sense fixed above (equal $\bar c$, equal span,
equal $K$; unequal rows). By the argument of
Section~\ref{sec:classfunction}(ii) this is structural, and it is therefore not
an independent comparison of two rules: they realise the same optimum.

\paragraph{The quota rule as configured is suboptimal in selection cost.}
The quota constraint prevents the greedy optimum from being reached, and the
recorded total selection cost of the quota pool exceeds that of the exhaustive
pool at every one of the $84$ points, never falling below it. The constraint is
thus active at every geometry, and the quota arm is a genuinely different
selection rather than a relabelling. This is a statement about the splits
listed in Section~\ref{sec:q4} only; a different admissible split
$(n_{\mathrm{s}},n_{\mathrm{q}})$ with $n_{\mathrm{s}}+n_{\mathrm{q}}=M$ could
in principle attain the unconstrained optimum, and no such split was tested.

\paragraph{A lower selection cost did not produce a lower $K$.}
Despite being worse under the objective that the ranking minimises, the quota
rule did not require a larger coupled space. Table~\ref{tab:dk} counts
$\Delta K$ of Eq.~\eqref{eq:dK} as (lower, equal, higher) over the geometries
of each block.

\begin{table}[htbp]
\centering
\begin{tabular}{llll}
\hline
system & cost & $M=n-\rsp$ & $M=n-\rsp-1$\\
\hline
$\mathrm{H_2O}$ & NC       & $(0,10,0)$, $\overline{\Delta K}=+0.00$ & $(0,10,0)$, $+0.00$\\
$\mathrm{H_2O}$ & variance & $(0,10,0)$, $+0.00$ & $(0,10,0)$, $+0.00$\\
$\mathrm{N_2}$  & NC       & $(5,4,2)$, $-0.27$ & $(5,5,1)$, $-0.45$\\
$\mathrm{N_2}$  & variance & $(1,10,0)$, $-0.09$ & $(1,10,0)$, $-0.09$\\
\hline
\end{tabular}
\caption{Quota rule against the exhaustive rule, counted as
(lower, equal, higher) in $\Delta K=K_{\mathrm{quota}}-K_{\mathrm{exhaustive}}$
over the $10$ ($\mathrm{H_2O}$) or $11$ ($\mathrm{N_2}$) geometries of each
block, with the mean $\overline{\Delta K}$. Negative favours the quota rule.}
\label{tab:dk}
\end{table}

The effect is confined to $\mathrm{N_2}$ and predominantly to the NC cost. For
$\mathrm{H_2O}$ the quota pool is more expensive in $\bar c$ at all $40$ points
($10$ geometries $\times$ $2$ cost functions $\times$ $2$ budgets) yet returns
exactly the same $K$ at all $40$; under the parity variance the $\mathrm{N_2}$
difference is a single geometry at each budget. Where a difference does appear
its sign is not consistent: over $\mathrm{N_2}$ and both cost functions the
range of $\Delta K$ is $[-2,+3]$ at $M=n-\rsp$ and $[-2,+1]$ at $M=n-\rsp-1$,
the largest single departure being $+3$ at $R=1.2$\,\AA.

This bears on the interpretation of the NC score. The exhaustive rule minimises
$\bar c$ exactly, whereas $K$ is measured after the orbital optimisation of
Eq.~\eqref{eq:objective}; the two are separated by that optimisation. A rule
that is worse in $\bar c$ and no worse in $K$ is therefore not a contradiction.
What these data establish is narrow and is stated as such: \emph{in the present
configuration} --- two molecules in a minimal basis, the quota splits listed
above, one orbital start --- ordering candidate rows by $\bar c$ did not order
the resulting $K$, and at $40$ of the points a strictly worse $\bar c$ produced
an identical $K$. Whether $\bar c$ is monotonically related to $K$ in general
is not addressed here, and no such claim is made; the observation is that the
relationship failed to hold in this configuration. In particular the exhaustive
column must not be read as a lower bound on $K$: it minimises the selection
cost, not the coupled dimension.

\subsection{What these data do not establish}
\label{sec:q4limits}

\begin{enumerate}
\item \textbf{No scatter estimate.} A single orbital start (the identity) was
run at each point, so no within-configuration variance is available. The
observed differences in $K$ are of the same order as the shifts seen between
budgets at fixed rule, and a difference of one or two units must be read as
``no detectable difference at one start''. Deciding whether the quota effect is
real requires repeating a subset with randomised starts at fixed pool; that
study is identified in Section~\ref{sec:conclusions} and has not been performed.
\item \textbf{$\mathrm{H_2O}$ is uninformative at the reduced budget.} At
$M=\dim\Qcal-1$ the decoupled energy already lies within $\varepsilon_E$ of
$E_{\mathrm{target}}$ at nine of ten $\mathrm{H_2O}$ geometries, giving $K=1$. These
are floor values, not agreement between rules, and they carry no discriminating
information; the $\mathrm{N_2}$ rows carry the entire Q4(ii) comparison.
\item \textbf{One quota split per budget.} The constrained rule was run at
$(n_{\mathrm{s}},n_{\mathrm{q}})=(2,2)/(3,3)$ and $(2,1)/(3,2)$ only. The
complementary splits were not tested, so ``the quota rule is suboptimal in
$\bar c$'' is established for these splits and not for quota selection as such.
\item \textbf{Spin-symmetric ground set.} Candidates are products of $\Zb_p$
(Section~\ref{sec:classfunction}(0)), so every optimum reported here is an
optimum over $\iota(\FF^{n})$, not over $\FF^{N}$. A spin-resolved search may
admit lower-cost pools at the same $M$.
\item \textbf{Scope.} Two molecules in a minimal basis, one orbital-rotation
parameterisation, and one reference convention. Nothing here establishes
behaviour at larger basis sets or for systems whose exact group is not a point
group supplemented by particle-number parities.
\item \textbf{Near-degenerate roots.}\begin{sloppypar}
The parent FCI is solved with \texttt{pyscf.fci.direct\_spin1} at
\texttt{conv\_tol}\,$=10^{-10}$ and \texttt{max\_cycle}\,$=10^{4}$, from
PySCF's deterministic initial guess; no
random seed enters. Where the FCI gap is small --- $\mathrm{H_2O}$ at
$R\ge2.33$\,\AA, gap $4.5\times10^{-4}$\,Ha --- no explicit tie-break is
imposed on near-degenerate eigenvectors, so the selected vector is reproducible
for a fixed build but is not guaranteed to be stable across BLAS or PySCF
versions. Energies are unaffected; overlap-ordered quantities in principle are.
\end{sloppypar}
\end{enumerate}


%=====================================================================
% New section. Place after Section~\ref{sec:q4} and before
% "Verification of the implementation".
% Requires the generated table gen_budget_tradeoff.tex.

\section{Budget and the block-dimension trade-off}
\label{sec:budget}

\subsection{Motivation}

Section~\ref{sec:classfunction} showed that at $M=\dim\Qcal$ every admissible
pool spans the whole quotient, so the sector partition cannot distinguish
selection rules. Two questions follow. First, whether the agreement between
cost-ranked rules reported in Section~\ref{sec:q4} survives when that
degeneracy is removed. Second, what is paid, in block dimension, for reducing
the budget. Both are addressed by repeating the campaign of
Section~\ref{sec:q4} at $M=\dim\Qcal-1$, that is $M=3$ for $\mathrm{H_2O}$ and $M=5$
for $\mathrm{N_2}$, with all other settings unchanged.

\subsection{The agreement is not an artefact of the budget}

At $M=\dim\Qcal-1$ the spanned subspace is a proper subspace of $\Qcal$, so
different rules \emph{may} select different subspaces. The exhaustive and
iterative rules nonetheless span the identical subspace at all $42$ points of
the reduced-budget campaign, and again return identical $K$. The agreement of
Section~\ref{sec:q4} is therefore a property of greedy optimality on the
matroid, as Proposition~\ref{prop:coset} and
Section~\ref{sec:classfunction}(ii) predict, and not a consequence of the
degenerate budget.

We record this explicitly because the converse inference is available and
wrong: observing agreement at $M=\dim\Qcal$ alone would not have distinguished
the structural explanation from the degeneracy, since at that budget the
partition is rule-independent by construction and the span test cannot fail.

\subsection{Cost of reducing the budget}

\begin{table}[htbp]
\centering
\begin{tabular}{llrrrr}
\hline
molecule & budget & $M$ & $\Dmax$ & retained dim & $K$ range\\
\hline
$\mathrm{H_2O}$ & $n-\rsp$ & 4 & 21 & 133 & 4--8\\
$\mathrm{H_2O}$ & $n-\rsp-1$ & 3 & 27,31 & 133 & 1--3\\
$\mathrm{N_2}$ & $n-\rsp$ & 6 & 120 & 1824 & 4--26\\
$\mathrm{N_2}$ & $n-\rsp-1$ & 5 & 156 & 1824 & 3--25\\
\hline
\end{tabular}

\caption{Effect of reducing the budget by one generator, at the exhaustive rule
with the NC cost. $\Dmax$ is the largest joint exact-and-approximate sector;
``retained dim'' is the total dimension of the sectors surviving the exact
filter; the $K$ range is over the geometries of the scan. Where two values of
$\Dmax$ appear, they occur at different geometries.}
\label{tab:tradeoff}
\end{table}

Three features of Table~\ref{tab:tradeoff} are relevant.

\paragraph{The retained dimension is invariant.}
The dimension of the retained exact sector, $\sum_{s}\dim(s)$ over sectors
surviving the exact filter, is unchanged between budgets. This is required:
that dimension is fixed by $\Ecal$ and the chosen exact sector alone, and the
LAS rows only subdivide it. Its invariance is used here as a consistency check
on the two campaigns rather than as a result.

\paragraph{Removing one generator coarsens the partition.}
\begin{sloppypar}
$\Dmax$, the largest joint exact-and-approximate sector and hence the largest
block that must be diagonalised, rises from $120$ to $156$ for $\mathrm{N_2}$, a
factor of $1.300$ at every geometry, and from $21$ to $27$ for $\mathrm{H_2O}$ at nine
of ten geometries, a factor of $1.286$. The exception is $\mathrm{H_2O}$ at
$R=2.5$\,\AA, where $\Dmax=31$ (factor $1.476$); this is also the only $\mathrm{H_2O}$
geometry at which the reduced budget does not reach the floor discussed below,
and the two observations have the same origin --- the leading sector there is
not large enough to absorb the target state. This increase is the direct cost
of the reduced budget.
\end{sloppypar}

\paragraph{The coupled dimension falls, but not monotonically.}
For $\mathrm{N_2}$ the coupled dimension decreases at eight of eleven geometries, the
range falling from $4$--$26$ to $3$--$25$, since a coarser partition places more
of the target state in the leading sector. It increases at the remaining three,
$R=1.7$--$1.9$\,\AA\ ($23\to25$, $18\to19$, $14\to15$). Because the partition at
fixed budget is identical across rules, an increase in $K$ at reduced budget
cannot originate in the partition; the only remaining channel is the orbital
optimisation converging to a different stationary point.

We do not attribute it further. In particular these three geometries are
\emph{not} those at which the quota rule departs most from the cost-ranked
rules: the largest $|\Delta K|$ between rules occurs at $R=1.2$\,\AA\ at
$M=n-\rsp$ and at $R=1.7$\,\AA\ at $M=n-\rsp-1$, and at $R=1.8$\,\AA\ the two
rules agree exactly at the reduced budget. The budget-induced and rule-induced
departures therefore do not localise to the same geometries, and there is no
evidence here that they are the same phenomenon. Both are candidates for the
multi-start study of Section~\ref{sec:conclusions}, which is the only way to
establish whether either exceeds the scatter of the optimiser.

\paragraph{$\mathrm{H_2O}$ reaches the floor.}
At $M=3$ the decoupled energy is already within $\varepsilon_E$ of
$E_{\mathrm{target}}$ at nine of ten geometries, so $K=1$ and no coupling is
required. No comparative information about selection rules is available from
these points, and they are excluded from the Q4(ii) reading for that reason.

\subsection{A consequence for budget selection}

For $\mathrm{H_2O}$ in this basis the reduced budget is not merely tolerable but
preferable on both axes that the procedure exposes. Dropping from $M=4$ to
$M=3$ removes one generator, and in exchange $K$ falls from the range $4$--$8$
to $1$ at nine of ten geometries and $3$ at the tenth, while $\Dmax$ rises from
$21$ to $27$ (to $31$ at $R=2.5$\,\AA). The coupled problem is thereby
eliminated at nine of ten geometries at a cost of one moderately larger block.

Whether that exchange is favourable depends on the cost model, and the present
data do not settle it: the coupled stage scales with $K$ and the number of
retained sectors, the decoupled stage with $\Dmax$ per block, and the two are
not reduced to a common measure anywhere in this report. What the data do
establish is that the trade exists and is not one-sided --- reducing the budget
is not simply a degradation --- and that the budget is therefore a parameter
worth scanning rather than fixing at the bound $M=n-\rsp$. For $\mathrm{N_2}$
the same reduction lowers $K$ at eight of eleven geometries and raises it at
three, so no comparable statement holds there.


%=====================================================================
\section{Verification of the implementation}
\label{sec:verification}

Three checks with independently known answers were performed.

\begin{enumerate}[leftmargin=1.8em,itemsep=2pt]
\item \textbf{Sector dimensions.} The dimension of the fully pinned exact
  sector was computed by direct enumeration of determinants at fixed
  $(N_\alpha,N_\beta)$ and compared with the dimension obtained from the
  Clifford tapering; the two agree for both systems and for every subgroup of
  the exact set.
\item \textbf{Quotient enumeration.} The number of distinct nonzero classes
  produced by the coset-representative routine was checked against
  $2^{\,N-r}-1$ at three register sizes, with exact agreement.
\item \textbf{Rotation convention.} Equation~\eqref{eq:convention} was
  determined empirically rather than assumed, by requiring the recorded values
  of Eq.~\eqref{eq:nc} to be reproduced; the transpose alternative is wrong by
  up to a factor of $78$. We note that the FCI energy is invariant under
  orbital rotation and therefore cannot validate the convention: only a
  rotation-sensitive quantity can.
\end{enumerate}

Two checks remain outstanding and are recorded as such: a brute-force
verification of Eq.~\eqref{eq:commutator} against explicit sparse construction
of $Z(g)$ and $[\Ht,Z(g)]$ for a small case, and an independent check of the
Givens parameterisation against a directly constructed rotation matrix.

%=====================================================================
\section{Reproducibility record}
\label{sec:record}

The following are recorded for all $126$ points; a representative point is
given in full ($\mathrm{N_2}$, $R=1.8000$\,\AA, irrep-block packing, iterative
rule, NC objective). Machine-readable records for every point are produced by
the script listed in Appendix~\ref{sec:artifacts}.

\begin{description}[leftmargin=1.6em,style=nextline,itemsep=3pt]

\item[Frame and exact set.]
Qubit ordering: interleaved JW, $2p\mapsto p\alpha$, $2p+1\mapsto p\beta$.
Rotation convention: Eq.~\eqref{eq:convention}, Givens parameterisation.
Allowed rotation group: intra-irrep blocks of $D_{2h}$.
Exact spatial rows as support masks $\{192,288,617\}$, i.e.\
$Q_{\pi x}=\{6,7\}$, $Q_{\pi y}=\{5,8\}$, $Q_u=\{0,3,5,6,9\}$; with
$P_\alpha,P_\beta$ this gives $r=5$.
Exact sector eigenvalue label $(1,1,0,0,0)$.

\item[Problem specification.]
$N=20$ qubits, $n=10$ spatial orbitals, $r=5$, $M=7$.
Candidate mode: polynomial current-frame, Eq.~\eqref{eq:Cpoly}, of size $55$ in
the spatial subregister.
Initial frame $F^{(0)}=I$ (rows $\Zb_p$); initial rotation $x^{(0)}=0$.
Completion rule: the ranked scan is continued to $N-r$ accepted rows, with no
separate completion routine.
Tie-break: ascending (cost, frame index).

\item[Per-iteration scan.]
$20$ macroiterations. The complete candidate list, its score, its sorted
position and the accept/reject reason are recorded for every candidate at every
macroiteration ($16$ entries at macroiteration $0$; $41$--$42$ thereafter).
Final LAS rows $\{118,481,733,528,822,937,777\}$; no auxiliary rows, since
$M=n-\rsp$ in the subregister.

\item[Binary spans.]
Intermediate and final spans recorded at all $20$ macroiterations, both for
$\Ecal\cup\Gcal$ and for the complete ranked basis $\Bcal$. Final span in RREF:
$\{1,2,4,8,16,32,64,128,256,512\}$, i.e.\ the full register.

\item[Orbitals, optimiser, state.]
Rotation parameters
$(-0.01702,\,-0.00033,\,-0.02624,\,0.00140,\,-0.23293,\,0.03532)$.
$C_A$ before $0.514304$, after $0.308498$.
Optimiser L-BFGS-B, exit on relative reduction of the objective below the
machine-precision threshold, step tolerance $10^{-5}$, gradient tolerance
$10^{-6}$, $87$ recorded objective evaluations.
Reference reconstructed at each trial rotation as the exact FCI ground state;
$E_{\mathrm{FCI}}=-107.483457$\,Ha.
Exact-symmetry leakage $1-W=4\times10^{-16}$.

\item[Decoupled energy and coupled dimension.]
$\Edec=-107.411543$\,Ha. The full sequence $\lambda_{\min}(H^{(K)})$ is
recorded, beginning
\[
-107.408879,\quad -107.439625,\quad -107.453531,\quad -107.469484\ \mathrm{Ha}.
\]
$500$ eigenstates per block, $\varepsilon_E=1.6\times10^{-3}$\,Ha.

\end{description}

%=====================================================================
\section{Conclusions}
\label{sec:conclusions}

\begin{enumerate}[leftmargin=1.8em,itemsep=4pt]

\item Exact-parity quotienting is valid only if the stored generators commute
  with the rotated Hamiltonian \emph{at the geometry in use}. Generators stored
  as fixed orbital-index sets fail this along a dissociation scan, silently,
  and the failure destroys the convergence of Eq.~\eqref{eq:K} rather than
  degrading it. More generally, identifying symmetry generators by orbital
  index is incompatible with PES scans of any kind: a level crossing
  redefines the operator without any change in the stored data. Generators must
  be re-derived from orbital symmetry labels at each geometry, or carried
  between geometries by an explicit continuation.

\item The retained reference weight $W$ of Eq.~\eqref{eq:W} is a sufficient and
  inexpensive run-time test for this failure. It costs one overlap sum, and
  $W<1$ makes Eq.~\eqref{eq:K} unsatisfiable irrespective of the pool. Deriving
  supports per geometry from the HF irreducible-representation labels restores
  $W=1$ at all $126$ points and with it the convergence of the criterion.

\item Unrestricted $\mathrm{SO}(n)$ rotation is inadmissible when the exact set
  contains point-group parities: a rotation mixing different irreducible
  representations carries the generator out of the diagonal Pauli group, so the
  Clifford reduction no longer applies. The measured leakage reaches $28.5\%$,
  and no compensating benefit is observed --- of $119$ paired points, $118$
  agree exactly in $K$ with the irrep-block packing.

\item The polynomial current-frame library of Eq.~\eqref{eq:Cpoly} recovers the
  exhaustive optimum, and this is guaranteed rather than fortuitous. By
  Proposition~\ref{prop:coset} the cost is constant on cosets of $\Ecal$, so it
  is well defined on $\Qcal$; the classes accepted by Eq.~\eqref{eq:ranktest}
  are the independent sets of a linear matroid over $\FF$; and greedy ranking
  by a well-defined weight on a matroid returns a minimum-weight independent
  set of size $M$ for every $M$. Two rules that both attain that optimum must
  therefore carry equal total weight and span the same subspace, differing at
  most by the choice of coset representative, which is unobservable. This is
  what is measured: over the $84$ paired points of
  Sections~\ref{sec:q4} and~\ref{sec:budget} the two rules agree in total cost
  to $2.7\times10^{-14}$, in span at every point, and in $K$ at every point,
  while never returning the same row set.

\item At the full budget the selection rule cannot affect the partition, and at
  the reduced budget it does not. With $M=n-\rsp$ the LAS span is the entire
  quotient for every rule, so the partition --- and hence the retained
  dimension, the block count and $\Dmax$ --- is identical by construction. The
  campaign was therefore repeated at $M=n-\rsp-1$, where the span is a proper
  subspace and a rule \emph{may} change the partition. It does not: the
  cost-ranked rules again span the same subspace and return the same $K$ at all
  $42$ points, confirming that the agreement of Conclusion~4 follows from
  matroid optimality and not from the degenerate budget.

  The constrained quota rule is the only rule that differs. It is more
  expensive in selection cost at all $84$ points, yet does not require a larger
  coupled space: for $\mathrm{H_2O}$ it returns exactly the same $K$ at all
  $40$ points, and for $\mathrm{N_2}$ under the NC cost it is lower at $5$ of
  $11$ geometries at each budget, with $\overline{\Delta K}=-0.27$ and $-0.45$
  and a sign that is not consistent. Ordering candidates by the selection cost
  therefore did not order $K$ in this configuration. Because only one orbital
  start was run, no scatter estimate exists and this difference is reported as
  an observation, not a significant effect.

\end{enumerate}

\paragraph{Work outstanding.}
Three items remain, in order of importance. First, a robustness study over
several orbital starts at fixed pool: the quota difference above, and the three
$\mathrm{N_2}$ geometries at which $K$ rises when the budget is reduced
(Section~\ref{sec:budget}), are both of a size that a single start cannot
separate from optimiser scatter. Second, extension of the candidate
representation to spin-resolved rows, so that the ground set is $\FF^{N}$
rather than $\iota(\FF^{n})$ (Section~\ref{sec:classfunction}(0)) and the
exhaustive control can be repeated over the full quotient. Third, recording the
selection costs at the converged orbitals rather than at the starting frame,
which affects the reported cost bookkeeping but not $K$. The two verification
checks flagged in Section~\ref{sec:verification} also remain to be performed.

%=====================================================================
\appendix
\section{Artifacts}
\label{sec:artifacts}
\small
\begin{tabular}{@{}p{0.40\textwidth}p{0.56\textwidth}@{}}
\toprule
Artifact & Producer \\
\midrule
per-geometry exact parity matrices &
  \texttt{scripts/write\_default\_exact\_parity.py -{}-chk} \\
endpoint grid &
  \texttt{scripts/trillium\_rerun\_fixed\_grid.sh} \\
$W$, exact-block count, Eq.~\eqref{eq:K} verification &
  \texttt{scripts/audit/audit\_06\_verify\_rerun.sh} \\
reproducibility record, per point &
  \texttt{scripts/audit/audit\_07\_repro\_record.sh} \\
exhaustive control &
  \texttt{scripts/exhaustive\_quotient\_check.py} \\
self-consistent exhaustive fixed point &
  \texttt{scripts/exhaustive\_selfconsistent.py} \\
software version manifest &
  emitted by the audit scripts into \texttt{tables/audit/} \\
\bottomrule
\end{tabular}
\normalsize

\end{document}
